\hypertarget{appendix-2c568f05ee34}{}\label{referencestart:appendix-2c568f05ee34}
\section{Latent-variable models and generative flows}\label{appendix-2c568f05ee34-latent-variable-models-and-generative-flows}

\textbf{Curriculum depth:} Developed · \textbf{First seen:} University courses · \textbf{Latest development:} University courses

\subsection{Definition and mechanism}\label{appendix-2c568f05ee34-definition-and-mechanism}

Generative models describe how data could be produced, often by first sampling a latent variable \(z\) and mapping it to an observation \(x\). A variational autoencoder (VAE) learns an encoder distribution \(q(z|x)\) and decoder \(p(x|z)\); its evidence lower bound balances reconstruction likelihood against a KL penalty that keeps the latent distribution tractable. A generative adversarial network (GAN) trains a generator against a discriminator, creating a game that can produce sharp samples but may be unstable or collapse to limited modes. Normalizing flows use invertible transformations so the data density can be computed with a change-of-variables Jacobian determinant. These methods solve different tradeoffs among likelihood, sample quality, latent structure, and optimization stability. Their assumptions matter: a VAE's posterior family may be too simple, a GAN score is not a calibrated density, and a flow's invertibility constrains architecture. Generated samples are not evidence that the model has learned the true causal process.

\subsection{Prerequisites}\label{appendix-2c568f05ee34-prerequisites}

\begin{itemize}
\tightlist
\item
  \hyperref[appendix-4815cd87ab8b]{Probability density and change of measure}
\item
  \hyperref[appendix-29f4bea4c2bb]{Constrained optimization and inverse kinematics}
\item
  \hyperref[appendix-65788587ef78]{Research validity and ablation}
\item
  \hyperref[appendix-40df0855b4a7]{Automatic differentiation and backpropagation}
\end{itemize}

\subsection{Used by}\label{appendix-2c568f05ee34-used-by}

\begin{itemize}
\tightlist
\item
  \href{../07-Vision-XR-and-Playable-Worlds/main.pdf\#nameddest=04-visual-computing-ii-ch031-generative-models-learn-distributions-and-samplers}{Generative models learn distributions and samplers (Vision, XR, and Playable Worlds, Vision, XR, and Playable Worlds, Chapter 2)}
\item
  \hyperref[appendix-c5b332e400e7]{Radiance fields and Gaussian splatting}
\end{itemize}

\subsection{Full teaching explanation}\label{appendix-2c568f05ee34-full-teaching-explanation}

\begin{itemize}
\tightlist
\item
  \href{../07-Vision-XR-and-Playable-Worlds/main.pdf\#nameddest=04-visual-computing-ii-ch031-generative-models-learn-distributions-and-samplers}{Generative models learn distributions and samplers (Vision, XR, and Playable Worlds, Vision, XR, and Playable Worlds, Chapter 2)}
\end{itemize}

\subsection{Recall and next depth}\label{appendix-2c568f05ee34-recall-and-next-depth}

\begin{itemize}
\tightlist
\item
  State the VAE reconstruction/KL tradeoff, GAN adversarial objective, and flow change-of-variables role.
\item
  Which assumption or training failure distinguishes each family?
\end{itemize}

\subsection{Source and navigation}\label{appendix-2c568f05ee34-source-and-navigation}

Teaching narrative: \href{../07-Vision-XR-and-Playable-Worlds/main.pdf\#nameddest=04-visual-computing-ii-ch031-generative-models-learn-distributions-and-samplers}{open the main explanation (Vision, XR, and Playable Worlds, Vision, XR, and Playable Worlds, Chapter 2)}. Browse the \hyperref[appendix-a6d556a4ef50-university-graphics-vision-and-photography---independent-series]{University graphics and vision index}. Additional chapter connections: \href{../02-Mathematics-for-Visual-Computing/main.pdf\#nameddest=03-visual-computing-i-ch012-numerical-integration-and-monte-carlo-estimation}{Numerical integration and Monte Carlo estimation (Mathematics for Visual Computing, Mathematics for Visual Computing, Chapter 16)}; \href{../05-Spatial-Worlds-Engine-State-and-Simulation/main.pdf\#nameddest=04-visual-computing-ii-ch029-optimization-constraints-and-inverse-kinematics}{Optimization, constraints, and inverse kinematics (Spatial Worlds, Engine State, and Simulation, Spatial Worlds, Engine State, and Simulation, Chapter 10)}; \href{../02-Mathematics-for-Visual-Computing/main.pdf\#nameddest=01-foundations-i-ch019-a-linear-equation-is-a-measurement-and-several-equations-are-geometry}{A linear equation is a measurement, and several equations are geometry (Mathematics for Visual Computing, Mathematics for Visual Computing, Chapter 3)}.

Caption/source records: \href{../07-Vision-XR-and-Playable-Worlds/main.pdf\#nameddest=04-visual-computing-ii-ch036-sc-13---stanford-cs231n-deep-learning-for-computer-vision-spring-2025-lecture-13-generative-models-1}{SC-13 (Vision, XR, and Playable Worlds, Vision, XR, and Playable Worlds, Chapter 5)}, \href{../07-Vision-XR-and-Playable-Worlds/main.pdf\#nameddest=04-visual-computing-ii-ch036-sc-14---stanford-cs231n-deep-learning-for-computer-vision-spring-2025-lecture-14-generative-models-2}{SC-14 (Vision, XR, and Playable Worlds, Vision, XR, and Playable Worlds, Chapter 5)}.
